% Laget automatisk av elevark/lag_elevark.py fra opplegg/relativistisk-bevegelsesmengde.html. Endringer her blir overskrevet.
\documentclass[11pt]{article}
\usepackage{elevark}
\begin{document}
\elevarkhode{Fysikk 2 · Spesiell relativitetsteori}{Relativistisk bevegelsesmengde og energi – arbeidsark}{https://lektorbo.no/sim/relativistisk-bevegelsesmengde.html}{Simuleringen}

\section*{Oppgaver}

Bruk Konstant kraft og lysfarten\footnote{\url{lektorbo.no/sim/relativistisk-bevegelsesmengde.html}}. Start med et proton og $F = 1{,}0 \cdot 10^{-12}\,\mathrm{N}$. $m_{\mathrm{p}} = 1{,}67 \cdot 10^{-27}\,\mathrm{kg}$ og $c = 3{,}00 \cdot 10^{8}\,\mathrm{m/s}$.

\begin{oppgaveliste}
\oppgave \fase{gjett}{Gjett}Regn ut klassisk hvor lang tid det tar før protonet når lysfarten. Hva tror du skjer i virkeligheten? Kjør simuleringen og se på fartsgrafen.
\svar
\oppgave \fase{test}{Regn}Når $t = mc/F$, er $p = mc$. Bruk $p = γmv$ til å vise at farten da er $v = c/\sqrt{2}$. Hva er $γ$? Sjekk med simuleringen.
\svar
\oppgave \fase{forklar}{Forklar}Den samme impulsen gir den samme bevegelsesmengden i begge beskrivelsene. Hvorfor blir farten likevel forskjellig?
\svar
\oppgave \fase{test}{Undersøk}Velg grafen for bevegelsesmengde mot fart. Hvor stor prosentvis feil gjør vi med $p = mv$ ved $0{,}10c$ og ved $0{,}50c$? Stemmer tommelfingerregelen om 10~\% av lysfarten?
\svar
\oppgave \fase{test}{Regn}Velg romskipet på 50 tonn med en motor som skyver med 1,0~MN. Hvor lang tid tar det før romskipet har $0{,}99c$?
\svar
\oppgave \fase{forklar}{Forklar}Velg grafen for kinetisk energi mot fart. Hvorfor trengs det mye mer energi for å gå fra $0{,}90c$ til $0{,}9991c$ enn fra $0{,}70c$ til $0{,}90c$? Hvorfor kan ingen partikkel med masse nå lysfarten?
\svar
\end{oppgaveliste}

\end{document}
